\section{Project Overview}

\subsection{Project Purpose, Scope, and Objectives}

Many patients, particularly people with chronic conditions or prescriptions
from several healthcare professionals, have difficulty determining whether
their medications are safe to use together. Risks include drug--drug
interactions, duplicate active ingredients, contraindications, missed doses,
and possible adverse drug reactions. Prescription readability and different
brand names for related active ingredients make the problem harder.

The project will produce a Medication Safety Assistant that helps users answer
the question: ``What potential safety risks should I be aware of in my current
medications, and when should I seek professional advice?'' The product supports
prescription-image OCR, manual entry, typed requests, and voice input.
Speech-to-text transcribes a user's voice, while NLP identifies relevant
intents and medication information from the resulting text. The system requires
the user to verify all OCR-extracted or voice-transcribed details before using
them in an assessment. It first identifies Vietnamese products
and active ingredients against the medicine-registration data published by the
\href{https://dichvucong.dav.gov.vn/congbothuoc/index}{Drug Administration of
Vietnam (DAV)}, then maps the normalized active
ingredients to DDInter 2.0 entities for international drug--drug interaction
and severity checking. AI may explain structured system findings,
but it must not independently determine medication safety.

The project objectives are:

\begin{itemize}
    \item Provide multimodal medication and request input through prescription
          OCR, manual entry, typed conversation, and voice-to-text, with an
          editable confirmation step before further processing;
    \item Apply NLP to verified text to identify user intent and relevant
          medication information without allowing the language model to make
          independent medication-safety conclusions;
    \item Identify Vietnamese
          medicine names, active ingredients, strengths, and registration
          information using data published by the Drug Administration of
          Vietnam;
    \item Map the identified active ingredients to corresponding DDInter 2.0
          entities and assess supported drug--drug interaction risks and their
          severity using international interaction data;
    \item Clearly distinguish detected risks, no known interaction within the
          available data, and incomplete or unassessed cases;
    \item Provide understandable safety guidance and route high-risk or
          uncertain cases to professional review;
    \item Create reminders only from confirmed prescription schedules and
          record adherence and possible adverse drug reactions; and
    \item Deliver a demonstrable Web prototype within the five-sprint course
          schedule.
\end{itemize}

The initial release excludes advanced timing analysis, drug--food interaction
checks, pill recognition, handwritten prescription support, and integration
with real healthcare services.

\subsection{Assumptions and Constraints}

\begin{longtable}{|L{0.18\textwidth}|L{0.35\textwidth}|L{0.35\textwidth}|}
\hline
\textbf{Type} & \textbf{Description} & \textbf{Effect on the project}\\ \hline
Assumption & Users can review and correct OCR results and voice transcripts before assessment. & Safety checks run only on user-verified medication details, never directly on an unconfirmed OCR result or speech transcript.\\ \hline
Assumption & The medicine-registration data published by the Drug Administration of Vietnam (DAV) is available as the primary source for identifying Vietnamese medicinal products and their active ingredients. & The team must preserve source provenance and account for corrections, changes, supplements, or withdrawals that may not be reflected in the original registration publication.\\ \hline
Assumption & Vietnamese active ingredients can be mapped reliably to corresponding DDInter 2.0 entities for interaction checking. & Mapping rules, ambiguous matches, coverage, and permitted DDInter access must be validated early; unresolved mappings remain unassessed rather than safe.\\ \hline
Assumption & Professional review is simulated for the course prototype. & The prototype demonstrates the workflow but does not connect to real clinicians or clinical systems.\\ \hline
Assumption & A responsive Web application is sufficient for desktop and mobile users. & The team can maintain one codebase and test on representative modern browsers.\\ \hline
Schedule constraint & The course requires three phases and five sprints with a small three-person team. & Scope must be prioritized around one end-to-end medication safety workflow.\\ \hline
Safety constraint & The system must not prescribe, stop, replace, or modify prescription medication. & All messages must use cautious language and direct users to qualified professionals when appropriate.\\ \hline
Data constraint & Missing data or no detected interaction is not proof of safety. & The interface must display coverage limitations and an incomplete/unassessed state.\\ \hline
Reminder constraint & New reminders for medication involved in a severe interaction require professional review and confirmation. & A safety gate is required before reminder activation; unrelated medication reminders may proceed.\\ \hline
Privacy constraint & Prescription images, medication lists, health information, and symptoms are sensitive. & The prototype must minimize collected data, restrict access, avoid exposing secrets, and use non-identifying test data where possible.\\ \hline
Technical constraint & Initial OCR supports printed or digital Vietnamese prescriptions, not handwriting. & Low-confidence and unsupported input must be corrected manually or reported as unresolved.\\ \hline
Technical constraint & Voice-to-text accuracy varies with accent, pronunciation, background noise, microphone quality, and medication terminology. & The interface must display an editable transcript and require confirmation; low-confidence or unresolved medication terms must not be assessed automatically.\\ \hline
Privacy constraint & Voice recordings and transcripts may contain sensitive health information. & Audio collection requires clear permission, minimal retention, restricted access, and deletion after transcription when storage is not necessary.\\ \hline
\end{longtable}

\subsection{Project Deliverables}

\begin{longtable}{|L{0.24\textwidth}|L{0.41\textwidth}|L{0.13\textwidth}|L{0.12\textwidth}|}
\hline
\textbf{Deliverable} & \textbf{Description / acceptance condition} & \textbf{Owner} & \textbf{Target}\\ \hline
PA1 planning package & Initial Project Plan, Vision Document, 3--6 UI screens, two weekly reports, and accessible tracking evidence. & Whole team & Sprint 1\\ \hline
Requirements baseline & Reviewed vision, prioritized backlog, multimodal-input use cases, non-functional requirements, and traceability to interview needs. & Long & Sprint 2\\ \hline
Architecture and UX baseline & System context, main workflow, data/interface decisions, and tested UI prototype covering OCR, manual, typed, and voice input. & Trang / Thành & Sprint 2\\ \hline
Multimodal intake proof of concept & OCR extraction, speech-to-text transcription, NLP intent/medication processing, editable confirmation, and consistent verified-text output. & Long / Trang / Thành & Sprint 3\\ \hline
Risk-checking proof of concept & Prepared DAV medicine data, reviewed active-ingredient mappings to DDInter, and a controlled drug--drug interaction query with explicit result states. & Long & Sprint 3\\ \hline
Integrated test build & End-to-end Web build covering multimodal medication intake, verification, normalization, safety results, and safety guidance. & Whole team & Sprint 4\\ \hline
Final prototype release & Stable course prototype including voice input, confirmed-schedule reminders, basic adherence/ADR recording, regression results, documentation, and demo material. & Whole team & Sprint 5\\ \hline
Project evidence repository & Versioned source, decisions, test evidence, sprint records, risk register, meeting notes, and weekly reports accessible to teaching staff. & Trang & Continuous\\ \hline
\end{longtable}

